% !TEX program = xelatex
\documentclass[11pt,letterpaper]{article}

\usepackage[top=1in, bottom=1in, left=1.15in, right=1.15in]{geometry}

\usepackage{fontspec}
\usepackage{unicode-math}

\setmainfont{Latin Modern Roman}[
  Extension=.otf,
  UprightFont=lmroman10-regular,
  BoldFont=lmroman10-bold,
  ItalicFont=lmroman10-italic,
  BoldItalicFont=lmroman10-bolditalic,
]
\setsansfont{Latin Modern Sans}[
  Extension=.otf,
  UprightFont=lmsans10-regular,
  BoldFont=lmsans10-bold,
  ItalicFont=lmsans10-oblique,
  BoldItalicFont=lmsans10-boldoblique,
]
\setmonofont{Latin Modern Mono}[
  Extension=.otf,
  UprightFont=lmmono10-regular,
  ItalicFont=lmmono10-italic,
  Scale=0.85,
]

\usepackage{xcolor}
\usepackage{titlesec}
\usepackage{enumitem}
\usepackage{fancyhdr}
\usepackage{hyperref}
\usepackage{xurl}
\usepackage{ragged2e}
\usepackage{microtype}

\usepackage{../ac-paper-essay}

\hypersetup{
  colorlinks=true,
  linkcolor=acpurple,
  urlcolor=acpurple,
  pdfauthor={@jeffrey},
  pdftitle={My People — a directory for Sage},
}

\pagestyle{fancy}
\fancyhf{}
\renewcommand{\headrulewidth}{0pt}
\fancyhead[L]{\footnotesize\color{acgray}\itshape internal — for sage}
\fancyhead[R]{\footnotesize\color{acgray}\thepage}

\newcommand{\ac}{Aesthetic\textcolor{acpink}{.}Computer}
\newcommand{\np}{notepat}

% per-person heading: name + role chip
\newcommand{\person}[2]{%
  \vspace{0.9em}\noindent%
  {\normalfont\bfseries\fontsize{12pt}{14pt}\selectfont\color{acdark}#1}%
  \hspace{0.6em}{\footnotesize\color{acgray}\itshape #2}%
  \par\vspace{0.15em}\nopagebreak%
}

\newcommand{\field}[1]{\textbf{\small\color{acdark}#1}\hspace{0.3em}}

\begin{document}

\thispagestyle{empty}
\vspace*{0.4in}

\essaytitle{my people}
\essaydate{a working directory for sage \textperiodcentered{} may 2026}

\noindent sage --- this is the list i promised, written down so we can
actually use it. everyone on it is connected to \ac{} in some load-bearing
way: they ship code with me, they make music with me, they host me, they
hire me, or they print me on the wall of an institution. each entry has
\textbf{what they do well}, \textbf{what we have already made together},
and \textbf{where i think they could plug into a product}. it is not
exhaustive --- it is the people i would actually call. if anyone is
missing, the gap is mine.

the document is internal. don't share past us without asking. emails and
phone numbers are deliberately left off where they would be awkward in
your hands; if a collaboration gets real i'll make the intro warm.

\acbreak

\section{inner circle}

\person{Fia Benitez \textit{(@fifi)}}{partner \textperiodcentered{} calarts \textperiodcentered{} inbox + opportunity curator}

\field{practice} project coordination, written voice, taste; runs the
opportunity inbox alongside me (\texttt{honeydo.txt}); forwards calls,
grants, residencies; catches things i wouldn't see.

\field{what we've made} the entire opportunity pipeline that funnels into
\ac{}'s grant + residency activity. her edits land on bios, drafts,
public-facing copy.

\field{product angle} any product with a written-voice surface (onboarding,
emails, the prompt's text) benefits from her ear. she is also the most
honest first reader i have --- if a product premise reads as bullshit to
fia it almost certainly is.

\acbreak

\section{ac core --- people who have actually shipped code}

these are the humans whose names appear in \texttt{git log} on the main
repo, or whose work is built so deeply into the project that the codebase
itself remembers them. ordered by depth of involvement, not vanity.

\person{Artur Matveichenkov}{ac dev-friend \textperiodcentered{} mostly imessage \textperiodcentered{} booted-by patron}

\field{practice} graphics + drawing tools, painting interfaces; the
\texttt{sign} demo he shipped out of puerto rico; the brush in \ac{}
that's queued to be rewritten in the kidlisp graphical editor.

\field{what we've made} ssh access to the dev infra; a recurring
``artur painting story'' thread in the project's \texttt{TODO.txt};
\texttt{JtoA} (jeffrey-to-artur), the speech-recognition prototype
meant to make us communicate more complex thoughts than imessage allows;
on the \texttt{booted-by} screen at the grey tier --- one of the people
who actually paid into the project. our ongoing back-channel is imessage,
not git --- ideas land there first.

\field{product angle} the highest-bandwidth peer i have on drawing /
brush / paint-engine work; the relationship that has already produced one
named pattern in the codebase (\texttt{JtoA}: building tools to talk to
each other better). if a product is shaped like a drawing instrument or a
two-person creative protocol, artur is the call before anyone else.

\person{Georgica Pettus \textit{(@georgica, gh: ggcajp)}}{ac core contributor}

\field{practice} javascript piece development, prompt engineering, LLM
experiments, generative work inside the AC piece API.

\field{what we've made} 130+ commits across the main repo. long-running
contributor; one of the people who has seen the system evolve from inside
the codebase.

\field{product angle} if a product needs someone who already speaks the
piece API natively, she is one of three people in the world. strong on
prompt + LLM glue work.

\person{Tina Tarighian \textit{(gh: ttarigh)}}{ac core contributor}

\field{practice} ac piece development, javascript / kidlisp, ui work.

\field{what we've made} 80+ commits across pieces and tooling; she is one
of the people i have actually onboarded into the codebase.

\field{product angle} solid pair for anything that turns a piece into a
user-facing app or installation. piece-API fluent.

\person{Esteban Uribe}{ac contributor \textperiodcentered{} swift / native}

\field{practice} performance optimization, keyboard input layers, swift /
mac native ui. recently shipped the global keyboard-disable path inside
Menu Band.

\field{what we've made} \texttt{slab/menuband/} performance work in 2026
--- about 80 commits across the menubar synth and adjacent native surfaces.

\field{product angle} the one to call if a product crosses the browser/native
boundary on mac. native input + audio scheduling specialist.

\person{Miles Peyton \textit{(gh: mileshiroo)}}{ac contributor}

\field{practice} ac piece development; the contributor i know least well
of this group, but the commits are real.

\field{product angle} unknown ceiling --- worth a coffee before any product
brief.

\person{Ida Pruitt \textit{(@ida, gh: idapruitt)}}{ac contributor \textperiodcentered{} early user}

\field{practice} piece-level javascript; the designated ``first user'' for
\texttt{niki.aesthetic.computer} testing in earlier eras of the project.

\field{what we've made} co-author of \texttt{sno} (with sage \& me); about
25 commits across pieces.

\field{product angle} the strongest user-research candidate i know who is
also fluent enough to ship the change herself.

\person{Maya Man}{ac contributor \textperiodcentered{} independent artist}

\field{practice} her own browser-native art practice; on \ac{} she
shipped the \texttt{@maya/sparkle} piece.

\field{product angle} a peer rather than a teammate --- but if a product
ever wants the web-as-medium voice she has, she'd be the call.

\person{Bash Elliott \textit{(@bash, gh: rackodo)}}{ac contributor \textperiodcentered{} user-host}

\field{practice} ac piece creation and hosting; the canonical example in
docs for \texttt{@handle/piece-name} resolution.

\field{product angle} more of a community-anchor than a co-builder.

\person{Dan / rcrdlbl \textit{(@ Tlon)}}{ac contributor \textperiodcentered{} backend / auth}

\field{practice} login / auth flows, identity surfaces. day job at Tlon
(the Urbit company), which is relevant for any product that wants
decentralized-identity rails.

\field{product angle} the one to call if a product touches identity,
login, or wants to think about a non-google-non-apple stack.

\acbreak

\section{studio peers \textperiodcentered{} falsework}

\person{Matt Doyle \& Yuehao Jiang}{co-principals \textperiodcentered{} falsework \textperiodcentered{} booted-by patrons (purple tier)}

\field{practice} run \textbf{Falsework}, the studio building \textit{Spider
Lily} (unreal engine 5.6 game project on perforce / helix core at
\texttt{falsework.helixcore.io}). matt also shows up in the whistlegraph
collaborator list at the Schneider Museum, so the relationship predates
falsework itself.

\field{what we've made} the \texttt{false.work/} directory in this repo
holds the unreal build pipeline (windows + mac + ios) for spider lily ---
the studio and \ac{} share build infrastructure. both of them appear
together on the \texttt{booted-by} screen at the purple tier; matt's
brother \textbf{Eric Doyle} sits at the brown tier alongside them.

\field{product angle} the only game-studio-shaped relationship i have.
unreal + perforce + mac/ios codesigning expertise that would otherwise be
weeks of bring-up. if a product wants any 3D-game-engine reach, or wants
to think about \ac{} as a tool that ships into another studio's
pipeline, this is the live example.

\acbreak

\section{music collaborators}

the \texttt{pop/} lane has its own peer set --- people i make tracks with,
or who run nearby on the same scene.

\person{Quargs Greene}{music collaborator \textperiodcentered{} they/them \textperiodcentered{} computer club peer}

\field{practice} solo electronic music; sample-based composition; releases
on their own publishing.

\field{what we've made} \textbf{50/50 writer/performer split} on a track
they are building on top of \texttt{trancepenta} for their 3rd solo album.
i sent stems on 2026-05-21. derivative is clearance-safe because the
trancepenta source SFX are CC0.

\field{product angle} an early signal that \texttt{pop/} stems are an
\textit{output} other artists want to use. if we ever package a stems-as-API
product, quargs is the proof-of-life user.

\person{Will \textit{(@c\_robo on instagram)}}{music tech \textperiodcentered{} fresh calarts grad \textperiodcentered{} chromebook orchestrator}

\field{practice} music tech; just graduated calarts. runs a many-chromebook
distributed-music rig --- the cheap-hardware-as-instrument lineage in a
calarts-pedagogy frame. orchestrates a roomful of them as a single
instrument.

\field{product angle} the cheapest possible network-of-devices music
product has a name and a face here. anything that wants to treat
underpowered hardware as an ensemble (rather than an apology) should be
sketched with will in the room first. obvious overlap with \texttt{pop/}'s
chromebook-tier render targets if we ever want a live-performance shape
for those tracks.

\acbreak

\section{curatorial peers}

\person{Celine Wong Katzman}{associate curator (creative time) \textperiodcentered{} sfpc board}

\field{practice} associate curator at \textbf{Creative Time}; sits on
the \textbf{Board of Directors} of the School for Poetic Computation
(alongside American Artist, Salome Asega, Tega Brain, Taylor Levy,
Che-Wei Wang, Galen Macdonald). curatorial reach + institutional read.

\field{what we've made} the relationship i trust enough to write
confidential research \textit{for}. the scott-moore private brief in
\texttt{papers/scott-moore-profile/} is explicitly prepared for celine
--- ``Private brief --- prepared for Celine, Not for publication.'' she
gets the analytical layer that doesn't go on the public platter.

\field{product angle} the curator i'd brief first on any product that
needs to be read as serious work in the contemporary-art conversation
--- and the one most likely to know whether an institutional pilot
(a creative time commission, an sfpc residency, a board placement) is a
plausible next step.

\acbreak

\section{ucla social software \textperiodcentered{} where i'm in residence}

\person{Casey Reas}{residency host \textperiodcentered{} ucla sosoft \textperiodcentered{} processing co-creator}

\field{practice} co-creator of Processing (with ben fry, at MIT media lab);
language-as-art practitioner; teaches at UCLA; cofounded UCLA Social Software
(``sosoft'') with lauren lee mccarthy.

\field{relationship} hosts me as \textbf{author in residence} at sosoft
in 2026. the canonical CV line is \textit{Author in Residence, UCLA Social
Software (Casey Reas)}. he is the host. mccarthy is the studio's cofounder,
not a co-host.

\field{product angle} the most credentialed person in this document for
any product that wants academic shelter or language-design pedigree.
\ac{}'s entire ``language as the artwork'' stance has a clean line back
to processing --- which is useful when explaining the project to people
who don't yet read the work.

\person{Lauren Lee McCarthy}{ucla sosoft cofounder}

\field{practice} public/social interface art (\textit{Lauren}, \textit{Someone},
\textit{Bot Party}); p5.js progenitor; teaches at UCLA.

\field{relationship} cofounded sosoft \textit{with} reas; she is the
co-founder of the studio i'm in residence at, \textbf{not} my co-host.
mention her only when the audience knows her.

\field{product angle} adjacent rather than active --- her practice and
mine share an audience but we haven't built anything together.

\acbreak

\section{client work \textperiodcentered{} ongoing professional}

these are the gigs that are paying for hardware, server time, and album
covers. listed because they are real relationships and real bandwidth --- if
a product needs a friendly first customer, one of these three is the door.

\person{Rebecca Bonneville, NP}{client \textperiodcentered{} the resilience institute}

\field{practice} functional / naturopathic psychiatry; depth psychotherapy;
neurosomatic and regenerative-medicine work. running a small private
practice she is reframing as a four-door institute.

\field{what we've made} a poc gateway page live under
\texttt{aesthetic.computer/bp-0cf886709d1753/} (cinematic, fludd
macrocosmus, EB garamond, ``the resilience institute''). full rebrand of
\texttt{resilientminds.health} in scope; alchemical / 16th--17th-c.
anatomical aesthetic; aman.com as the experiential north star.

\field{product angle} the most concrete ``serious wellness brand built on
\ac{} infra'' example we have. if we ever sell a small-clinic /
practitioner product, she is the prototype customer.

\person{Loretta Darling Nicola}{client \textperiodcentered{} nevada natural living}

\field{practice} wellness products (copper sulfate, ADEK, supplements);
sells under the ``dr. vk for life'' brand.

\field{what we've made} \texttt{drvkforlife.com} on shopify, custom DNS
through cloudflare, marketing PDFs + product info sheets. invoice
AC-DRVK-2026-0519-001, \$1{,}550, net 30.

\field{product angle} less of a creative-product fit; relevant if a future
product wants a wellness-commerce ramp.

\person{Thomas Lawson}{client \textperiodcentered{} artist site}

\field{practice} painter / writer with a long-running site at
\texttt{thomaslawson.com}.

\field{what we've made} site maintenance; the deploy path goes godaddy
sftp into a mu-plugins setup --- not glamorous, but real income and a real
artist's site.

\field{product angle} weak product fit. listed for completeness.

\acbreak

\section{service partners \textperiodcentered{} humans, not vendors}

\person{Solveig Fernlund}{architect \textperiodcentered{} brooklyn}

\field{practice} residential interior architecture; pre-war renovation.
i've worked with her on the kapin/jacobs kitchen renovation in brooklyn
(164 mac donough st).

\field{product angle} not software-shaped. but if a product ever needs a
physical-space partner who understands the difference between a kitchen
that works and a kitchen that performs, she's it.

\person{@mollysoda}{piece collaborator \textperiodcentered{} \texttt{selfie}}

\field{practice} internet-native artist; she co-designed the decorated
photographic selfie palette behind \ac{}'s \texttt{selfie} piece.

\field{product angle} a clean example of \ac{}'s piece-as-collaboration
pattern; could be re-activated if a product wants a known internet-art
voice on a specific piece.

\acbreak

\section{institutional footprint --- not collaborators}

worth knowing about, but not on the call sheet. listing here so you have
the full picture.

\begin{itemize}[leftmargin=1.4em]
  \item \textbf{KADIST} (san francisco) and \textbf{SMK \textperiodcentered{} the
    national gallery of denmark} (copenhagen) hold ac work in their
    collections. they are \textit{collectors}, not collaborators --- the
    distinction matters. they signal seriousness, not bandwidth.
  \item \textbf{Tezos Foundation} --- ecosystem grant submitted 2026-05-22
    for the kidlisp keeps v12 work; awaiting Q2 review. not a person, but
    a check the project might cash.
  \item \textbf{LACMA} --- 2026 grant application in process, biennial
    structure (prototype 2027, premiere 2028). same shape: an institution,
    not a collaborator.
\end{itemize}

\acbreak

\section{where this leaves us}

sage --- the people who actually ship with me are a small group: you, me,
\textbf{artur}, \textbf{georgica}, \textbf{tina}, \textbf{esteban},
\textbf{ida} (lighter). the music lane has \textbf{quargs} as the first
external user of a \ac{} output and \textbf{will} (\@c\_robo) as the
chromebook-orchestra angle if we ever want a live-performance shape.
\textbf{matt + yuehao} (Falsework) are the live example of \ac{} infra
living inside another studio's pipeline. the curatorial read comes from
\textbf{celine}; the academic cover from \textbf{reas}. the gigs keep the
lights on. everyone else is a known node, not an active edge.

a product we build together would almost certainly be small at the start
--- the two of us, maybe artur on drawing-tools work, or one of the ac
core contributors when piece-API fluency is required. the music-collab
pattern (stems-as-product), the piece-as-collaboration pattern
(\texttt{selfie} with mollysoda), and the build-pipeline-as-shared-infra
pattern (\texttt{false.work/} with falsework) are the three shapes that
have already proven they ride out into the world.

let's pick a shape and a name. i'll keep this list current.

\colophon{written for sage jenson, may 2026. internal; not for distribution.}

\end{document}
